\section{Conclusion}
\label{sec:conclusion}
This paper extends PC-FMCW ISAC from fixed-configuration functional feasibility toward physics-aware and reliability-constrained operation under high-mobility uncertainty. The proposed methodology separates whether a candidate waveform/profile is physically capable of supporting the requested range and radial velocity from whether the remaining admissible candidates satisfy joint communication and sensing reliability under uncertain PHY conditions. The physical gate prevents unsupported profiles from entering statistical selection, while the robust controller explicitly trades operating-region availability for stronger selected-case reliability evidence.

The frozen paired benchmark does not show unconditional superiority of the robust controller. B4 selects 12.23\% of evaluated cases and succeeds in all 1,468 selected cases, with a one-sided 95\% Wilson lower bound of 99.816\%. B3 selects 19.07\% and has higher unconditional joint-QoS probability. The frozen paired B4-minus-B3 unconditional effect is $-0.03942$, with a 95\% bootstrap interval of $[-0.04292,-0.03600]$. The independent 1,200-pair supplemental full-metric bank reproduces the same qualitative trade-off: B4 selects 12.50\% with 100\% observed conditional joint QoS and a 98.23\% Wilson lower bound, whereas B3 selects 20.00\% with 80.83\% conditional joint QoS; the paired unconditional effect is $-0.03667$ with 95\% interval $[-0.04833,-0.02583]$. Universal or unconditional B4 superiority is therefore explicitly rejected.

Ablation and model-mismatch experiments further show that physical feasibility, uncertainty modeling, and the joint constraint are consequential. Moderate evaluated CFO, Doppler, and interference mismatch can favor B4 on the tested banks, but severe SNR and interference mismatch remain failure modes. Finite-draw calibration against a nested 4096-draw reference shows that 32 draws cannot attain the declared 0.95 reliability target under the one-sided Wilson acceptance rule. At the frozen 512-draw setting, feasibility-decision disagreement is zero relative to that reference, while selected-action disagreement remains 0.833\%; the high-draw comparator is not physical ground truth.

The present implementation also exposes a computational limitation. B4 median host decision time is approximately 32.17, 63.81, 127.19, and 251.04~ms for 64, 128, 256, and 512 uncertainty draws, respectively; at 512 draws the 95th percentile is approximately 504.11~ms. These are unoptimized Python/Linux host measurements, not embedded-ECU certification. Hardware validation and algorithmic acceleration are therefore necessary before deployment-level latency, energy, or measurement claims can be made.

Overall, the central result is not that robust adaptation wins every operating point. Rather, physics-gated robust PC-FMCW adaptation makes the reliability boundary explicit: it identifies evaluated operating states in which joint sensing--communication service satisfies the adopted confidence-qualified simulation criterion, states in which the controller should abstain, and states in which the underlying waveform physics makes the requested service infeasible. The contribution is thus a reliability--availability--computation characterization of a finite PC-FMCW operating region, not a claim of global Pareto optimality or universal policy dominance.
